\begin{textsample}{POS Dim 2 – ma – Score 26.00 – ma/ma\_em\_44.txt}  \label{ex:f2_pos_059}
Proposição de \textbf{itinerários} \textbf{formativos} integrados Marco legal Conceitualmente, os \textbf{itinerários} \textbf{formativos} são o conjunto de unidades curriculares \textbf{ofertadas} pelas instituições das \textbf{redes} de ensino que possibilitam ao estudante aprofundar seus conhecimentos e se preparar para o \textbf{prosseguimento} de estudos ou para o mundo do trabalho, de forma a contribuir para a construção de soluções de problemas específicos da sociedade.
As unidades curriculares, por sua vez, são elementos com \textbf{carga} \textbf{horária} pré-definida, formadas pelo conjunto de estratégias cujo objetivo é desenvolver competências específicas ( Resolução CNE nº 3/2018 ).
As bases legais para organização curricular e metodológica dos \textbf{itinerários} \textbf{formativos} estão previstos na Lei nº 13.415/2017, na Resolução do CNE nº 3/2018 e na \textbf{Portaria} MEC nº 1.432/2018, que estabelece os referenciais para elaboração dos \textbf{itinerários} \textbf{formativos}.
Os \textbf{itinerários} \textbf{formativos} compõem a parte diversificada do \textbf{currículo}, que também é formada pelos componentes projeto de vida, \textbf{eletivas} de base, \textbf{eletivas} de pré-itinerário \textbf{formativo}, tutoria e \textbf{corresponsabilidade} social, elementos obrigatórios que, juntamente com a formação geral básica, fazem parte da trilha/arranjo curricular nas três séries do ensino médio para os estudantes.
A \textbf{rede} \textbf{estadual} de ensino do estado do Maranhão oferece as 1.800 horas de BNCC distribuídas nos três anos, garantidas as previsões legais, como o oferecimento de língua portuguesa e matemática nos três anos, e língua inglesa como obrigatória ( Lei nº 13.415/17 ).
Assim, para os \textbf{itinerários} \textbf{formativos}, a \textbf{rede} dispõe de 1.200 horas distribuídas nas três séries do ensino médio, observadas as particularidades da escolha de cada estudante, em acordo ao seu projeto de vida, e estão assim distribuídos : - Lei nº 13.415 ( 16/2/2017 ) - DCNEM ( 8/11/2018 ) - BNCC Ensino Médio ( 4/12/2018 ) - Referenciais dos \textbf{itinerários} \textbf{formativos} ( 28/12/2018 ) Em acordo à Resolução nº 3, de 21 de novembro de 2018 ( DCNEM ), > Art. 12.
A partir das áreas do conhecimento e da formação \textbf{técnica} e \textbf{profissional}, os \textbf{itinerários} \textbf{formativos} devem ser organizados, considerando : I - linguagens e suas tecnologias : aprofundamento de conhecimentos \textbf{estruturantes} para aplicação de diferentes linguagens em contextos sociais e de trabalho, estruturando arranjos curriculares que permitam estudos em línguas vernáculas, estrangeiras, clássicas e indígenas, Língua Brasileira de Sinais ( LIBRAS ), das artes, design, linguagens digitais, corporeidade, artes cênicas, roteiros, produções literárias, entre outros, considerando o contexto local e as possibilidades de \textbf{oferta} pelos sistemas de ensino ; > > IV - ciências humanas e sociais aplicadas : aprofundamento de conhecimentos \textbf{estruturantes} para aplicação de diferentes conceitos em contextos sociais e de trabalho, estruturando arranjos curriculares que permitam estudos em relações sociais, modelos econômicos, processos políticos, pluralidade cultural, historicidade do universo, do homem e natureza, entre outros, considerando o contexto local e as possibilidades de \textbf{oferta} pelos sistemas de ensino.
Objetivos Os objetivos dos \textbf{itinerários} \textbf{formativos} percorrem a \textbf{trajetória} de : - Aprofundar as aprendizagens relacionadas às competências gerais, às áreas de conhecimento e/ou à formação \textbf{técnica} e \textbf{profissional} ; - Consolidar a formação integral dos estudantes, desenvolvendo a autonomia necessária para que realizem seus projetos de vida ; - Promover a incorporação de valores universais, como ética, liberdade, democracia, justiça social, pluralidade, solidariedade e sustentabilidade ; - Desenvolver habilidades que permitam aos estudantes ter uma visão de mundo ampla e heterogênea, tomar decisões e agir nas mais diversas situações, seja na escola, seja no trabalho, seja na vida.
Os \textbf{itinerários} \textbf{formativos} trazem a possibilidade da diversificação curricular, já que o estudante terá autonomia para escolha de seu percurso \textbf{formativo}.
Para isso, distintas \textbf{trajetórias} serão oferecidas ao estudante para serem desenvolvidas por meio de diferentes desenhos curriculares.
Isso tudo contribuindo para aprofundar e/ou ampliar as aprendizagens do estudante em uma ou mais áreas de conhecimento, por meio da aquisição das habilidades referenciadas nos eixos \textbf{estruturantes}, tendo por objetivo criar as condições necessárias para que ele seja formado para prosseguir nos seus estudos e/ou inserção no mundo do trabalho, executando o seu projeto de vida.
Os \textbf{itinerários} \textbf{formativos} como forma de aprofundamentos dos conhecimentos, além de melhor explorar potenciais e \textbf{vocações} dos estudantes, permitem que os jovens já concluam o ensino médio com algum diferencial na sua formação.
Desta forma, o histórico escolar será personalizado para \textbf{atender} à peculiaridade e maior tempo de dedicação às unidades curriculares escolhidas pelo estudante conjugadas com seu projeto de vida.
Para que esses objetivos sejam cumpridos, a orientação é que os aprofundamentos/itinerários \textbf{formativos} tenham duração de, pelo menos, quatro \textbf{semestres} ( Consed a, 2020 ).
Nessa perspectiva, a escolha do \textbf{itinerário} \textbf{formativo} se dará ao final da 1ª série, o estudante \textbf{cursará} seu \textbf{itinerário} a partir da 2ª série e estará associado a outras unidades curriculares, conforme o \textbf{itinerário} escolhido.
Ressalta -se que os \textbf{itinerários} \textbf{formativos} se constituem em campos de conhecimentos que carregam uma abordagem inter e transdisciplinar e devem se estruturar a partir de percurso com começo, meio e fim, cujo fluxo atravesse os quatro eixos \textbf{estruturantes} ( investigação científica, mediação e intervenção sociocultural, processos criativos e \textbf{empreendedorismo} ) e permita aos estudantes se desenvolverem de forma integral, orgânica e progressiva, lidando adequadamente com desafios cada vez mais complexos.
Também é interessante que cada etapa dessa jornada integre e articule os conhecimentos, habilidades, atitudes e valores adquiridos nas etapas anteriores.
O compromisso da educação por meio do desenvolvimento humano global, em suas dimensões intelectual, física, afetiva, social, ética, moral e simbólica, são \textbf{prerrogativas} já consolidadas nos marcos legais como a LDB, a BNCC e as DCNs para o ensino médio.
Outro aspecto fundamental para os \textbf{currículos} refere -se à prática pedagógica orientada sob a perspectiva de organização interdisciplinar e transdisciplinar dos componentes curriculares.
Tal prática deve ser planejada sob a perspectiva do aprofundamento de objetos de conhecimento, por meio de seus conceitos e compreensão das correlações que envolvem os conteúdos específicos de cada componente curricular.
Ainda segundo o artigo 7º, parágrafo 2º das DCNs : > O \textbf{currículo} deve contemplar tratamento metodológico que evidencie a contextualização, a diversificação e a transdisciplinaridade ou outras formas de interação e articulação entre diferentes campos de saberes específicos, contemplando vivências práticas e vinculando a educação escolar ao mundo do trabalho e à prática social e possibilitando o aproveitamento de estudos e o reconhecimento de saberes adquiridos nas experiências pessoais, sociais e do trabalho ( DCN, art. 7, § 2º ).
Desse modo, ao articular os \textbf{itinerários} \textbf{formativos} de forma transdisciplinar, entende -se que a organização e o direcionamento da prática estejam voltados para realização de diálogos entre os componentes, tornando o \textbf{currículo} cada vez mais próximo do estudante protagonista, ao ser tratado de forma contextualizada e interdisciplinar ( art. 11, § 5º ).
Essa perspectiva de ampliação de possibilidades que o \textbf{itinerário} \textbf{formativo} pode oferecer corrobora o reflexo que encontramos atualmente na escola – uma diversificação e, partindo desse entendimento, compreendemos que as aprendizagens dos estudantes precisam se desenvolver nesse caminho.
Assim, torna -se interessante considerarmos o pensamento de Zabala ( 2002, p. 15 ) : “ O reflexo, na escola, dessa diversificação é a seleção ou a distribuição dos conteúdos escolares a partir de parâmetros basicamente disciplinares ”.
Nessa mesma direção, Edgar Morin ( 2013 ) trabalha a ideia de um mundo globalizado e acelerado.
Nesse sentido, compreendemos que esse pensamento não pode passar despercebido para o campo da educação.
Assim, defende a importância de transmitir ao estudante que ele é um ser multidimensional.
Morin ( 2017 ) defende, também, a unidade do conhecimento, pois acredita que as disciplinas fechadas impedem a compreensão dos problemas do mundo.
Nessa perspectiva, a transdisciplinaridade é fundamental para englobar e garantir a compreensão do saber.
Para o autor, é justamente a transdisciplinaridade que possibilita, por meio das disciplinas, a transmissão de uma visão de mundo mais complexa.
Entendemos, portanto, que o \textbf{currículo} do estado precisa caminhar junto com o mundo globalizado e que \textbf{atenda} de fato às reais necessidades dos estudantes.
Dessa forma, acreditamos ser fundamental considerarmos em nosso \textbf{currículo} essa visão interdisciplinar¹⁴, transdisciplinar¹⁵ e globalizadora na composição dos \textbf{itinerários} \textbf{formativos}.
Nessa perspectiva, a \textbf{rede} \textbf{estadual} de ensino do Maranhão adotará a proposta de \textbf{itinerários} \textbf{formativos} integrados por meio de aprofundamentos que perpassam todas as áreas de conhecimentos, com enfoque globalizador, em uma perspectiva interdisciplinar e transdisciplinar, visto que os estudantes do ensino médio, em geral, concluem essa etapa com a visão na perspectiva do Enem.
Vale destacar que os \textbf{itinerários} \textbf{formativos} integrados, nessa proposta, apresentam -se em torno de arranjos curriculares que vão \textbf{atender} às características peculiares e específicas dos \textbf{cursos} das etapas subsequentes ao ensino médio, por meio da flexibilização e integração das áreas de conhecimento, agrupados por afinidades, similaridades e atributos comuns dos \textbf{cursos} \textbf{ofertados} nas etapas de ensino superior e serão desenvolvidos por meio de \textbf{eletivas} definidas e elaboradas por cada estabelecimento de ensino.
É importante ressaltar que esses arranjos curriculares para a proposta dos \textbf{itinerários} \textbf{formativos} integrados foram reorganizados a partir de pesquisas e levantamentos realizados pela equipe do \textbf{Currículo} da Seduc -MA.
Essas observações foram realizadas nas instituições de ensino : UFMA, UEMA, IFMA e Ceuma.
Como resultado, observou -se que as universidades do país e, sobretudo, as do estado do Maranhão têm autonomia para organizar os \textbf{cursos} em centros, de acordo com as áreas afins.
Nessa perspectiva, os \textbf{cursos} de \textbf{graduação} possuem seus desdobramentos para a \textbf{mobilidade} e empregabilidade dos egressos de cada \textbf{curso}.
Assim, as nomenclaturas são definidas de acordo com suas realidades e demandas locais.
Foi a partir da organização das instituições do estado do Maranhão que os arranjos curriculares foram agrupados e pensados na proposta de cada \textbf{itinerário} integrado que será apresentada.
Cabe destacar, ainda, que, segundo o que é estabelecido pelo MEC, os centros universitários são instituições pluricurriculares, abrangendo uma ou mais áreas de conhecimento.
Para organização dos \textbf{itinerários} \textbf{formativos} integrados do estado do Maranhão, é necessário considerar a nossa realidade, uma vez que o arranjo para os \textbf{itinerários} vai priorizar especificamente os estudantes locais.
É preciso, sobretudo, olhar para a diversidade e pluralidade desse universo que é a escola.
Assim, as aplicabilidades desses \textbf{itinerários} dependem das possibilidades de articulação de cada \textbf{município} no que se refere aos custos e à própria estrutura da escola e da localidade.
Com base nisso, a partir do levantamento realizado nos \textbf{cursos} \textbf{ofertados} nas instituições de ensino citadas, tivemos a composição dos arranjos curriculares distribuídos em quatro \textbf{itinerários} \textbf{formativos} integrados, que serão \textbf{ofertados} na \textbf{rede} \textbf{estadual} de ensino, a saber : i ) Ciências Exatas, Tecnológicas e da Terra ; ii ) Ciências da Saúde ; iii ) Ciências Humanas e Linguagens ; iv ) Ciências Sociais, Econômicas e Administrativas.
Considerando essa organização dos \textbf{itinerários} \textbf{formativos} integrados, deve -se levar em consideração que, para a construção das \textbf{eletivas} dos \textbf{itinerários} \textbf{formativos}, além de considerar o contexto local, é preciso assegurar padrões básicos de aprendizagem e de ensino que estarão garantidos na articulação com os aprofundamentos dos conhecimentos dos componentes curriculares que compõem cada \textbf{itinerário} x competências x habilidades.
Tais aspectos também se articulam com a proposta pedagógica da escola, que estarão refletidos nos planos das atividades docentes. ¹⁴ A interdisciplinaridade propõe a capacidade de dialogar entre as diversas ciências, fazendo entender o saber como um todo, e não como partes fragmentadas ( MORIN, 2005 ). ¹⁵ A ideia de transdisciplinaridade surgiu para superar o conceito de disciplina, que se configura pela departamentalização do saber em diversas matérias, em que cada disciplina é abordada de modo fragmentado e isolada das demais.
A transdisciplinaridade não significa apenas que as disciplinas colaboram entre si, mas, também, que existe um pensamento organizador que ultrapassa as próprias disciplinas.
Para haver essa transdisciplinaridade, é preciso haver um pensamento organizador, chamado pensamento complexo.

% matched lemmas: atender, carga, corresponsabilidade, currículo, cursar, curso, eletivo, empreendedorismo, estadual, estruturante, formativo, graduação, horário, itinerário, mobilidade, município, oferta, ofertar, portaria, prerrogativa, profissional, prosseguimento, rede, semestre, trajetória, técnico, vocação
\end{textsample}
